\documentclass{amsart}
\def\publname{NOT A REAL JOURNAL}

\usepackage{enumitem}
\usepackage{hyperref}
\usepackage{subcaption}
\usepackage{amssymb}
\captionsetup{compatibility=false}

\usepackage{amsrefs}

%% \usepackage{subfigure}

\newcommand{\thing}[1]{$\mathcal{#1}$}

\def\startignorepylatexenc{}
\def\endignorepylatexenc{}


\def\sillything{\begin{proof}
    this is something that shouldn't really exist,
    but it $\alpha$ would \textit{pose a problem}.
\end{proof}}

\def%
%
%
\problem%
%
%
{\textit{yeah...} this is no good, right?}

%% metadata

\title[shorter title]{The title or whatever}

\author[The name that will be repeated over and over\dots.]{Full name one}
\address[sure]{The address or whatever}
\curraddr[something]{The current address}
\urladdr[url]{url addr}
\email{some\_email@gmail.com}

\author[name two]{Full name two}
\address[]{another address}
\curraddr[something]{second curr address}
\urladdr[url]{https://ams.org/}
\email{cak@ams.org}

\thanks{So this is where we say thanks to grant No.\ 325349M\@.}

\thanks{L.\ E.\ is the corresopnding author.}

\thanks{The first author gratefully acknolwedgaes the fruitful conversations they had with Jocaib Gewald.}
\thanks{The second author would like to thank the University of California Irvine for it's generous accomodations for research and
  spaces for exchanging ideas and collaborating productively.}

\subjclass{Primary 1586: Secondary 314, 159, 2653, 58579}

\keywords{Something, else, that, we, need}

\dedicatory{In memory of Vladimir Zelovsky}
\commby{Arlene Belfans}
\translator{Julia Gusteau}
\copyrightinfo{XXXX}{American Mathematical Society}

\begin{document}
\begin{abstract}[does this take an optional argument?]
  This is an abstract with some words. But no words that are particularly meaningful at all.\footnotemark But if you just use a straight \protect\footnote{stuff that won't appear anywhere at all: mingus} that won't work (ERROR). Well not compile erorr, but the mark will appear and the footnote contents will \dots dissappear\dots.
\end{abstract}\footnotetext{I would guess that this behaves like footnotes in captions?}

\makeatletter
\def\@maketitle{%
  \normalfont\normalsize
  \@adminfootnotes
  \@mkboth{\@nx\shortauthors}{\@nx\shorttitle}%
  \global\topskip42\p@\relax % 5.5pc   "   "   "     "     "
  \@settitle
  \ifx\@empty\authors \else \@setauthors \fi
  \ifx\@empty\@dedicatory
  \else
    \baselineskip18\p@
    \vtop{\centering{\footnotesize\itshape\@dedicatory\@@par}%
      \global\dimen@i\prevdepth}\prevdepth\dimen@i
    \fi
    \ifx\@empty\@commby
    \else
    \baselineskip18\p@
    \vtop{\centering{\footnotesize\@commby\@@par}%
      \global\dimen@i\prevdepth}\prevdepth\dimen@i
    \fi
  \@setabstract
  \normalsize
  \if@titlepage
    \newpage
  \else
    \dimen@34\p@ \advance\dimen@-\baselineskip
    \vskip\dimen@\relax
  \fi
} % end \@maketitle
\makeatother

\maketitle

Here is some text and a \problem. I think.

\section[but running head?]{A section title\protect\footnote{section one}}

\begin{figure}
  figure stuff
  \caption[]{A caption for a figure. This can't have a footnote in it, I believe\footnotemark}
\end{figure}\footnotetext{And that's where you're wrong. (If you give the {\tt \string\caption} command an empty optional argument---{\tt\string\caption[]}---so that the toc figurelist name is empty so that there's no issue with {\tt \string\par} in the caption body, then you can put {\tt \string\footnotemark} in the caption argument and then put {\tt \string\footnotetext} outside of the figure so that it behaves as normal: wild! But if you just use {\tt\string\footnote} in the caption argument even with giving it an empty optional argument, it will appear just like a {\tt\string\footnotemark} without a {\tt\string\footnotetext}.)}

Oh so does the footnote actually work?\footnote{Let's see\dots.} It looks like if you give\dots I don't know what I was originally saying in this sentence, but\dots.\footnote{This is a footnote that is very closely followed by another footnote.} Here, I'll show you:\footnote{Here's the other footnote, just as described\dots}.

This is an absolute mess of a document, but here's an equation

\[
\begin{aligned}
  X^2 &= \prod x\times iy -\frac{ \Theta \Omega}{\omega \hat{i}\bar{i+iz}}\\
  &= \sum_{i\in\infty}^{\frac{\eta\aleph}{\beth}}
  \end{aligned}
\]

which is a bit of a palette cleanser, maybe\footnote{is it?}.

% subfigure with \subfigure
%% \begin{figure}
%%   \subfigure[][Here's a caption\footnotemark]{subfigure one \qquad\qquad\qquad}
%%   \subfigure[][Here's the second caption\footnotemark]{subfigure two \qquad\qquad\qquad}
%%   \caption[]{Two subfigures\footnotemark}
%% \end{figure}
%% \footnotetext[3]{what will this go to?}
%% \footnotetext[4]{and of this?}
%% \footnotetext{Same thing in figure with subfigures.}

\vspace{12pt}

By default, you can put almost anything in a footnote, right?\footnote{This is maybe true---let's see:

\[
\prod_{e \in \mathbb{E}}^{\Omega} x^2 - \varphi^3
\]

well? Yes, it looks like everything is fine. Other than maybe\dots an enumerate?

\begin{enumerate}
\item first
  \item second
\end{enumerate} Yeah\dots no problem.
\begin{itemize}
\item bullet
\item list
\item of
\item stuff
\end{itemize} But you can't put a figure in a footnote, right?

\begin{center}
  \begin{tabular}{|c|c|}
  \hline
  stuff & more \\
  \hline
  And things & you hear?\\
  \hline
  \end{tabular}
\end{center}

But you cannot put a float in a footnote. Thank goodness.

You cannot put a verbatim environment either it seems.
{\tt\string\controlsequence}
} So it seems. Blech. 

%% subfigure with \begin{subfigure} %lookslike we can only load one of subcaption and subfigure
\begin{figure}
  \begin{subfigure}[t]{14.5pc}
    subfigure one
    \caption{Here's the first subcaption. It is relatievley long for a subfigure caption. Wait so \footnote{won't work without optional argument to caption, right?} wow. Well this must be specifically provided by subcaption, of course.}
  \end{subfigure}
  \begin{subfigure}[t]{14.5pc}
    subfigure two
    \caption{Here's the second subcaption\footnote{So you're telling me that subcaption supports footntoes in subcaptions?} Okay\dots wow you are.\footnote[26]{This is crazy\dots.}. Just absolutely bonkers!\footnote{Yes, it is.}}
  \end{subfigure}
  \caption{Two subfigures}
\end{figure}


\newpage

And finally some text on a new page!

\begin{bibdiv}
  \begin{biblist}
\bibitem[B]{brewster} S. Brewster, \textit{Automorphisms of the cohomology ring of finite Grassmann manifolds}, Thesis (Ph.D.)--The Ohio State University, ProQuest LLC, Ann Arbor, MI (1978), 102 pp. \url{http://gateway.proquest.com/openurl?url_ver=Z39.88-2004&rft_val_fmt=info:ofi/fmt:kev:mtx:dissertation&res_dat=xri:pqdiss&rft_dat=xri:pqdiss:7908118}
	

\bib{Aptekarev}{article}{
   author={Aptekarev, A.I.},
   title={Multiple orthogonal polynomials},
   journal={J. Comput. Appl. Math.},
   volume={99},
   year={1998},
   pages={423--447},
}

\bibitem[BH]{brewster-homer} S. Brewster and W. Homer, \textit{Rational automorphisms of Grassmann manifolds}, Proc. Amer. Math. Soc. \textbf{88} (1983), no. 1, 181--183, \url{https://doi.org/10.2307/2045137}.
	

\bib{Baker}{book}{
    AUTHOR = {Baker, G.A.},
    AUTHOR = {Graves-Morris, P.},
     TITLE = {Pad\'e{} approximants},
    SERIES = {Encyclopedia of Mathematics and its Applications},
    VOLUME = {59},
   EDITION = {Second},
 PUBLISHER = {Cambridge University Press, Cambridge},
      YEAR = {1996},
     PAGES = {xiv+746},
}


\bibitem[Do]{dold} A. Dold, \textit{Erzeugende der Thomschen Algebra }$\mathfrak{N}$, Math. Z. \textbf{65} (1956), 25--35, \url{https://doi.org/10.1007/BF01473868}.


\bib{BCVA}{article}{
   author={Beckermann, B.},
   author={Coussement, J.},
   author={Van Assche, W.},
   title={Multiple Wilson and Jacobi-Pi\~neiro polynomials},
   journal={J. Approx. Theory},
   volume={132},
   date={2005},
   number={2},
   pages={155--181}
}


\bibitem[Du]{duan} H. Duan, \textit{Self-maps of the Grassmannian of complex structures}, Compositio Math. \textbf{132} (2002), no. 2, 159--175, \url{https://doi.org/10.1023/A:1015885227445}.
	

\bib{Bultheel}{book}{
    AUTHOR = {Bultheel, A.},
     TITLE = {Laurent series and their {P}ad\'e{} approximations},
    SERIES = {Operator Theory: Advances and Applications},
    VOLUME = {27},
 PUBLISHER = {Birkh\"auser Verlag, Basel},
      YEAR = {1987},
     PAGES = {xii+274},
}


\bibitem[DF]{duan-fang} H. Duan and L. Fang, \textit{Homology rigidity of Grassmannians}, Acta Math. Sci. Ser. B \textbf{29} (2009), no. 3, 697--704, \url{https://doi.org/10.1016/S0252-9602(09)60065-5}.


\bib{Chihara}{book}{
   author={Chihara, T.S.},
   title={An Introduction to Orthogonal Polynomials},
   isbn={9780486479293},
   series={Mathematics and Its Applications},
   volume={13},
   publisher={Gordon and Breach Science Publishers, Inc.},
   year={1978},
}


\bibitem[DZ]{duan-zhao} H. Duan and X. Zhao, \textit{The classification of cohomology endomorphisms of certain flag manifolds}, Pacific J. Math. \textbf{192} (2000), no. 1, 93--102, \url{https://doi.org/10.2140/pjm.2000.192.93}.

\bib{MOPUC2}{article}{
   author={Cruz-Barroso, R.},
   author={D{\'i}az Mendoza, C.},
   author={Orive, R.},
   title={Multiple orthogonal polynomials on the unit circle. Normality and recurrence relations},
   journal={J. Comput. Appl. Math.},
   volume={284},
   year={2015},
   pages={115--132},
}



\bibitem[GH1]{glover-homer} H. Glover and B. Homer, \textit{Endomorphisms of the cohomology ring of finite Grassmann manifolds}, Geometric applications of homotopy theory (Proc. Conf., Evanston, Ill., 1977), I, pp. 170--193.

%\bib{DaeKui}{article}{
%    AUTHOR = {Daems, E.},
%    AUTHOR = {Kuijlaars, A.B.J.},
%     TITLE = {A {C}hristoffel-{D}arboux formula for multiple orthogonal
%              polynomials},
%   JOURNAL = {J. Approx. Theory},
%  FJOURNAL = {Journal of Approximation Theory},
%    VOLUME = {130},
%      YEAR = {2004},
%    NUMBER = {2},
%     PAGES = {190--202},
%      ISSN = {0021-9045,1096-0430},
   %MRCLASS = {42C05},
  %MRNUMBER = {2100703},
%MRREVIEWER = {K.\ Pan},
       %DOI = {10.1016/j.jat.2004.07.003},
       %URL = {https://doi.org/10.1016/j.jat.2004.07.003},
%}

\bib{DerSim18}{article}{
    AUTHOR = {Derevyagin, M.},
    AUTHOR = {Simanek, B.},
     TITLE = {Asymptotics for polynomials orthogonal in an indefinite
              metric},
   JOURNAL = {J. Math. Anal. Appl.},
  FJOURNAL = {Journal of Mathematical Analysis and Applications},
    VOLUME = {460},
      YEAR = {2018},
    NUMBER = {2},
     PAGES = {777--793},
}


\bibitem[GH2]{glover-homer coin} H. Glover and W. Homer, \textit{Fixed points on flag manifolds}, Pacific J. Math. \textbf{101} (1982), no. 2, 303--306, \url{http://projecteuclid.org/euclid.pjm/1102724776}.
	

\bib{GeronimusRelations}{article}{
    AUTHOR = {Geronimus, Ya.L.},
     TITLE = {On the trigonometric moment problem},
   JOURNAL = {Ann. of Math. (2)},
  FJOURNAL = {Annals of Mathematics. Second Series},
    VOLUME = {47},
      YEAR = {1946},
     PAGES = {742--761},
}

\bibitem[GS]{goswami-sarkar} A. Goswami and S. Sarkar, \textit{Endomorphisms of the Cohomology Algebra of Certain Homogeneous Spaces}, \url{https://arxiv.org/abs/2509.09363}


\bib{GeronimusBook}{book}{
    AUTHOR = {Geronimus, Ya.L.},
     TITLE = {Teoriya ortogonal\cprime nyh mnogo\v clenov},
 PUBLISHER = {Gosudarstv. Izdat. Tehn.-Teor. Lit., Moscow-Leningrad},
      YEAR = {1950},
     PAGES = {164},
}


\bibitem[Ho1]{hoffman} M. Hoffman, \textit{Endomorphisms of the cohomology of complex Grassmannians}, Trans. Amer. Math. Soc. \textbf{281} (1984), no. 2, 745--760, \url{https://doi.org/10.2307/2000083}.

\bibitem[Ho2]{hoffman-noncoin} M. Hoffman, Noncoincidence index of manifolds, Pacific J. Math. \textbf{115}(2) (1984), 373--383, \url{http://projecteuclid.org/euclid.pjm/1102708254}.


\bib{HueMan}{article}{
    AUTHOR={Huertas, E.J.},
    AUTHOR={Ma\~{n}as, M.},
     TITLE = {Mixed-type multiple orthogonal {L}aurent polynomials on the
              unit circle},
   JOURNAL = {J. Comput. Appl. Math.},
  FJOURNAL = {Journal of Computational and Applied Mathematics},
    VOLUME = {475},
      YEAR = {2026},
     PAGES = {Paper No. 117037},
}


\bib{Ismail}{book}{
   author={Ismail, M.E.H.},
   title={Classical and Quantum Orthogonal
Polynomials in One Variable},
   isbn={9780521782012},
   series={Encyclopedia of Mathematics and its Applications},
   Volume={98},
   publisher={Cambridge University Press},
   year={2005},
}

\bib{JNT}{article}{
    author={Jones, W.B.},
    author={Nj{\r a}stadt, O.},
    author={Thron, W.J.},
    title={Moment Theory, Orthogonal Polynomials, Quadrature, and Continued Fractions Associated with the unit Circle},
    journal={Bulletin of the London Mathematical Society},
    volume={21},
    number={2},
    year={1989},
    pages={113--152}
}

\bib{KNik}{article}{
    AUTHOR={Kozhan, R.},
    TITLE={Nikishin Systems on the Unit Circle},
  journal={},
   volume={},
   date={},
   number={},
   pages={under submission, arXiv:2410.20813},
}


\bibitem[HH]{hoffman-homer} M. Hoffman and W. Homer, \textit{On cohomology automorphisms of complex flag manifolds}, Proc. Amer. Math. Soc. \textbf{91} (1984), no. 4, 643--648, \url{https://doi.org/10.2307/2044817}.

    
\bibitem[L]{lin} X. Z. Lin, \textit{Geometric realization of Adams maps}, Acta Math. Sin. \textbf{27} (2011), no. 5, 863--870, \url{https://doi.org/10.1007/s10114-011-0164-y}.

\bibitem[M]{mandal} M. Mandal, Cohomology of generalized Dold manifolds, Thesis (Ph.D.), Homi Bhabha National Institute, The Institute of Mathematical Sciences (2024).
	
\bibitem[MS1]{mandal-sankaran} M. Mandal and P. Sankaran, \textit{Cohomology of generalized Dold spaces}, Topology Appl. \textbf{310} (2022), Paper No. 108040, 16 pp, \url{https://doi.org/10.1016/j.topol.2022.108040}.


\bib{KVMLOPUC1}{article}{
    AUTHOR={Kozhan, R.},
    AUTHOR={Vaktn{\"a}s, M.},
    TITLE={Angelesco and AT Systems on the Unit Circle},
    journal={},
    volume={},
    date={},
    number={},
    pages={under submission, arXiv:2410.12094},
}

\bib{KVChristoffel}{article}{
   author={Kozhan, R.},
   author={Vaktn\"{a}s, M.},
   title={Christoffel transform and multiple orthogonal polynomials},
      JOURNAL = {J. Comput. Appl. Math.},
    VOLUME = {476},
      YEAR = {2026},
     PAGES = {Paper No. 117121},
}

\bib{KVMOPUC}{article}{
    AUTHOR={Kozhan, R.},
    AUTHOR={Vaktn{\"a}s, M.},
    TITLE={Szeg\H{o} recurrence for multiple orthogonal polynomials on the unit circle},
    JOURNAL={Proc. Amer. Math. Soc.},
    VOLUME={152},
    NUMBER={11},
    YEAR={2024},
    PAGES={2983-2997},
    ISSN={1088-6826,0002-9939},
}

\bibitem[MS2]{mandal-sankaran2} M. Mandal and P. Sankaran, \textit{Cohomology and K-theory of generalized Dold manifolds fibred by complex flag manifolds}, \url{https://arxiv.org/abs/2407.03932}.

\bib{KVInterlacing}{article}{
    AUTHOR={Kozhan, R.},
    AUTHOR={Vaktn{\"a}s, M.},
    TITLE={Zeros of multiple orthogonal polynomials: location and interlacing},
    journal={Bull. of London Math. Soc.},
    volume={},
    date={},
    number={},
    pages={to appear, arXiv:2503.15122},
}


 
\bib{Kui}{article}{
   AUTHOR = {Kuijlaars, A.B.J.},
     TITLE = {Multiple orthogonal polynomial ensembles},
    JOURNAL = {Recent trends in orthogonal polynomials and approximation
              theory, Contemp. Math., Amer. Math. Soc., Providence, RI},
    VOLUME = {507},
     PAGES = {155--176},
      YEAR = {2010},
      ISBN = {978-0-8218-4803-6},
}

\bib{Applications}{article}{
   author={Mart{\'i}nez-Finkelshtein, A.},
   author={Van Assche, W.},
   title={WHAT IS...A Multiple Orthogonal Polynomial?},
   journal={Not. Am. Math. Soc.},
   volume={63},
   year={2016},
   pages={1029--1031},
}

\bib{MOPUC1}{article}{
   author={M{\'i}nguez Ceniceros, J.},
   author={Van Assche, W.},
   title={Multiple orthogonal polynomials on the unit circle},
   journal={Constr. Approx.},
   volume={28},
   year={2008},
   pages={173--197},
}

\bibitem[NS]{nath-sankaran} A. Nath and P. Sankaran, \textit{On generalized Dold manifolds}, Osaka J. Math. \textbf{56} (2019), no. 1, 75--90, Errata, \textbf{57} (2020), no. 2, 505--506, \url{https://projecteuclid.org/euclid.ojm/1547607627}.

\bib{Nikishin}{book}{
    AUTHOR = {Nikishin, E.M.},
    AUTHOR = {Sorokin, V.N.},
     TITLE = {Rational approximations and orthogonality},
    SERIES = {Translations of Mathematical Monographs},
    VOLUME = {92},
      NOTE = {Translated from the Russian by Ralph P. Boas},
 PUBLISHER = {American Mathematical Society, Providence, RI},
      YEAR = {1991},
     PAGES = {viii+221},
      ISBN = {0-8218-4545-4},
}

\bib{PeherstorferSteinbauer}{article}{
    author={Peherstorfer, F.},
    author={Steinbauer, R.},
    title={Characterization of orthogonal polynomials with respect to a functional},
    journal={Journal of Computational and Applied Mathematics},
    volume={65},
    year={1995},
    pages={339-355},
}


\bibitem[O]{O} L.S. O'Neill, \textit{On the fixed point property for Grassmann manifolds}, Thesis (Ph.D.)--The Ohio State University ProQuest LLC, Ann Arbor, MI (1974). 52 pp, \url{http://gateway.proquest.com/openurl?url_ver=Z39.88-2004&rft_val_fmt=info:ofi/fmt:kev:mtx:dissertation&res_dat=xri:pqdiss&rft_dat=xri:pqdiss:7511411}.

\bib{OPUC1}{book}{
   author={Simon, B.},
   title={Orthogonal Polynomials on the Unit Circle, Part 1: Classical Theory},
   isbn={0-8218-3446-0},
   series={Colloquium Lectures},
   Volume={54},
   publisher={American Mathematical Society},
   year={2004},
}

\bibitem[P]{Papadima} S. Papadima, \textit{Rigidity properties of compact Lie groups modulo maximal tori}, Math. Ann. \textbf{275} (1986), no. 4, 637--652, \url{https://doi.org/10.1007/BF01459142}.

\bib{OPUC2}{book}{
   author={Simon, B.},
   title={Orthogonal Polynomials on the Unit Circle, Part 2: Spectral Theory},
   isbn={978-0-8218-4864-7},
   series={Colloquium Lectures},
   Volume={54},
   publisher={American Mathematical Society},
   year={2005},
}


\bib{SimonL2}{book}{
   author={Simon, B.},
   title={Szeg\H{o}'s Theorem and Its Descendants: Spectral Theory for $L^2$ Perturbations of Orthogonal Polynomials},
   isbn={9780691147048},
   series={Porter Lectures},
   publisher={Princeton University Press},
   year={2011},
}

\bibitem[ST1]{shiga-tezuka} H. Shiga and M. Tezuka, \textit{Rational fibrations, homogeneous spaces with positive Euler characteristics and Jacobians}, Ann. Inst. Fourier (Grenoble) \textbf{37} (1987), no. 1, 81--106, \url{https://doi.org/10.5802/aif.1078}.

\bib{SzegoBook}{book}{
    AUTHOR = {Szeg\H o, G.},
     TITLE = {Orthogonal polynomials},
    SERIES = {American Mathematical Society Colloquium Publications},
    VOLUME = {Vol. XXIII, 4th ed},
   %EDITION = {Fourth},
 PUBLISHER = {American Mathematical Society, Providence, RI},
      YEAR = {1975},
     PAGES = {xiii+432},
}

  \bibitem[ST2]{shiga-tezuka2} H. Shiga and M. Tezuka, \textit{Cohomology automorphisms of some homogeneous spaces}, Singapore topology conference (Singapore, 1985), Topology Appl. \textbf{25} (1987), no. 2, 143--150, Errata, \textbf{34} (1990), no. 2, 207, \url{https://doi.org/10.1016/0166-8641(87)90007-1}.

\bib{Szego}{article}{
    author={Szeg\H{o}, G.},
    title={{\"U}ber die Entwickelung einer analytischen Funktion nach den Polynomen eines Orthogonalsystems},
    issn={0025-5831},
    journal={Mathematische Annalen},
    volume={82},
    year={1921},
    pages={188-212},
}

\bibitem[W]{wong} P. Wong, \textit{Fixed point theory for homogeneous spaces. II}, Fund. Math. \textbf{186} (2005), no. 2, 161--175, \url{https://doi.org/10.4064/fm186-2-4}.
	
\bib{NNRR}{article}{
   author={Van Assche, W.},
   title={Nearest neighbor recurrence relations for multiple
orthogonal polynomials},
   journal={J. Approx. Theory},
   volume={163},
   year={2011},
   pages={1427--1448},
}

\end{biblist}
\end{bibdiv}

\end{document}
